%  LaTeX support: latex@mdpi.com
%  manuscriptR1V2.tex: Round-1 revision rewrite.
%  Based on the MDPI Mathematics template from first_round/.
%  Internal bibliography (MDPI ACS format): references are embedded below;
%  biblio.bib is the editable source. To refresh after citation changes:
%  temporarily use \bibliography{biblio}, run pdflatex+bibtex, and copy the .bbl.

%=================================================================
\documentclass[mathematics,article,submit,moreauthors]{Definitions/mdpi}

%=================================================================
% MDPI internal commands - do not modify
\firstpage{1}
\makeatletter
\setcounter{page}{\@firstpage}
\makeatother
\pubvolume{1}
\issuenum{1}
\articlenumber{0}
\pubyear{2026}
\copyrightyear{2026}
%\externaleditor{Firstname Lastname}
\datereceived{ }
\daterevised{ }
\dateaccepted{ }
\datepublished{ }

%=================================================================
\usepackage{xcolor}

%=================================================================
% Full title of the paper (Capitalized)

\Title{Aggregation Semantics for Prioritization in Cooperative Automated Decision Making}

% Author Orchid ID: enter ID or remove command
\newcommand{\orcidauthorA}{0000-0003-3267-6753}
\newcommand{\orcidauthorB}{0009-0002-0183-7905}
\newcommand{\orcidauthorC}{0000-0002-7653-1452}

% Authors, for the paper (add full first names)
\Author{Pavel Novoa-Hern\'andez$^{1,*}$\orcidA{}, Mariia Godz$^{2}$\orcidA{} and David A. Pelta$^{2}$\orcidC{}}

% MDPI internal command: Authors, for metadata in PDF
\AuthorNames{Pavel Novoa-Hern\'andez, Mariia Godz and David A. Pelta}

% Affiliations / Addresses
\address{%
$^{1}$ \quad Dept. Computer Engineering and Systems, Universidad de La Laguna; pnovoahe@ull.edu.es\\
$^{2}$ \quad Dept. Computer Science and Artificial Intelligence, Universidad de Granada; \{mariiagodz,dpelta\}@ugr.es}

% Contact information of the corresponding author
\corres{Correspondence: pnovoahe@ull.edu.es}

% Abstract placeholder (to be rewritten with the new structure)
\abstract{To be rewritten.}

% Keywords
\keyword{
cooperative automated decision making;
aggregation operators;
multi-criteria decision making;
context-aware decision support;
prioritization
}

%%%%%%%%%%%%%%%%%%%%%%%%%%%%%%%%%%%%%%%%%%
\begin{document}

% ==============================================================================
% NOTA GLOBAL DE ESTILO:
% [Aborda 2.8 y M1.7]: Todo el documento debe ser reescrito en tercera persona
% (evitar "We provide", "We compare") y con una revisi\'on exhaustiva de ingl\'es.
% ==============================================================================

\section{Introduction}

% R1-1, R2-2
% Posicionar claramente la novedad.
% No proponemos un nuevo operador ni un nuevo método fuzzy MCDM.
% La contribución es metodológica: estudiar la agregación como mecanismo
% de gobernanza dentro de una arquitectura cooperativa ADM.

% R2-3
% Introducir explícitamente los research questions (RQ1-RQ4).

% R2-4
% Al final de la sección incluir la tabla cronológica de literatura
% (o referenciarla si se coloca en Background).

% R1-12, R2-6
% Justificar desde el principio por qué se estudian tres regímenes:
% compensatorio, parcialmente compensatorio y no compensatorio.

\subsection{Research Questions and Contributions}

% R2-3
% Formular RQ1-RQ4.
% Enumerar contribuciones metodológicas y experimentales.

%%%%%%%%%%%%%%%%%%%%%%%%%%%%%%%%%%%%%%%%%%%%%%%%%%%%%%%%%%%%

\section{Background and Related Work}

% R1-1, R2-2, R2-3, R2-4
% Construir el gap de investigación.
% Distinguir explícitamente el trabajo de:
% - fuzzy MCDM,
% - aggregation theory,
% - context-aware DSS,
% - cooperative ADM.

\subsection{Context-Aware and Cooperative Decision Support Systems}
\label{subsec:context-aware-dss}
% R2-2, R2-3
% Revisar Teng, Yao, Cronholm, Merkert, etc.
% Mostrar que el foco suele estar en reglas, ontologías o interacción.

The increasing adoption of Artificial Intelligence (AI) in decision-support systems has highlighted the need to complement predictive performance with contextual, organizational, and domain-specific knowledge. In many real-world settings, decisions cannot be based solely on statistical predictions because institutional policies, expert knowledge, regulatory requirements, and operational constraints must also be considered. As a result, a growing body of research has focused on context-aware and cooperative decision-support systems capable of integrating heterogeneous sources of information into the decision-making process.

One important research direction concerns the incorporation of explicit knowledge representations alongside data-driven models. Early examples include cognitive-agent architectures that combine symbolic rules with machine-learning mechanisms to support context-aware decision making \citep{TengTan2008}. Similarly, ontology-based clinical decision-support systems have employed contextual ontologies and clinical pathways to select context-appropriate recommendations for individual cases \citep{YaoKumar2012}. More recent developments have extended these ideas through semantic models, knowledge graphs, and hybrid AI architectures that combine predictive models with structured domain knowledge to improve decision quality and interpretability \citep{AbuRasheed2024,Casado2025Governance}.

A second research stream emphasizes cooperation between human and artificial intelligence. Rather than viewing decision support as a purely predictive task, these approaches consider decision making as a collaborative process in which machine-learning models provide evidence while human experts contribute contextual understanding, oversight, and governance. Cronholm and G\"obel \citep{CronholmGoebel2022}, for example, argue that effective decision support requires appropriate mechanisms for combining human and artificial intelligence rather than simply improving predictive accuracy. Likewise, recent surveys on human-centred AI highlight transparency, accountability, contextual awareness, and governance as fundamental requirements for trustworthy decision-support systems \citep{Fanfani2025Human}.

Related ideas also appear in cooperative and distributed decision-making environments. Multi-agent systems, collective-intelligence frameworks, and cooperative AI architectures frequently rely on multiple decision sources that must be reconciled before a final recommendation is produced \citep{Smirnov2021Methodology}. In these settings, contextual information often plays a coordinating role by constraining, guiding, or validating the outputs generated by predictive components. Similar concerns arise in governance-oriented decision-support systems, where contextual information may encode operational procedures, organizational priorities, safety requirements, or regulatory obligations \citep{Marot2022Perspectives,Casado2025Governance}.

Taken together, these studies clearly demonstrate the importance of
contextual knowledge in decision-support systems. However, their primary
focus typically lies in knowledge representation, ontology design,
human--AI interaction, governance mechanisms, or system architecture.
Once predictive evidence and contextual assessments have been generated,
the numerical integration process that combines both sources of
information is often treated as a secondary implementation detail.
Consequently, relatively little attention has been devoted to
understanding how alternative integration mechanisms affect the final
recommendations produced by the system.

This limitation becomes particularly visible in architectures that
explicitly separate predictive evidence from contextual assessment.
Recent work on Cooperative Automated Decision-Making Systems (CADEMAS)
\citep{NovoaHernandez2026} and its machine-learning instantiation
CADEMAS-ML \citep{Godz2026} exemplifies this design philosophy.
CADEMAS organizes heterogeneous automated decision-making components and
contextual knowledge into distinct layers, allowing predictive evidence
and contextual assessments to be generated independently before being
combined into a final recommendation. Such architectures make the
integration mechanism itself an explicit design choice rather than an
implicit implementation detail. As a result, the aggregation operator
determines how predictive evidence interacts with contextual
requirements, organizational priorities, or policy constraints, thereby
shaping the semantics of the resulting decision process.

This observation suggests that aggregation should not be viewed merely
as a numerical post-processing step. Instead, it can be interpreted as a
governance mechanism that regulates the interaction between predictive
and contextual information. Understanding the implications of different
aggregation semantics therefore becomes essential for the design of
cooperative and context-aware decision-support systems. The next
subsection reviews the aggregation literature from this perspective. 

\subsection{Aggregation Operators and Decision Semantics}
\label{subsec:aggregation-operators-and-decision-semantics}

% R1-1, R1-12, R2-2, R2-6
% Revisar OWA, t-norms, geometric operators, consensus aggregation.
% Introducir la idea de compensación.

Aggregation operators constitute one of the central building blocks of
multi-criteria decision making (MCDM), fuzzy decision support, and
collective decision-making systems. Their purpose is to combine multiple
sources of information into a single assessment while preserving, to a
greater or lesser extent, the decision semantics associated with the
original criteria. Consequently, aggregation has been studied
extensively from both theoretical and applied perspectives.

Among the most influential aggregation frameworks are Ordered Weighted
Averaging (OWA) operators \citep{Yager1988,Yager2005}, t-norm and
t-conorm based operators, geometric aggregation functions, and
non-additive integrals such as the Choquet and Sugeno integrals.
Collectively, these approaches provide a rich mathematical vocabulary
for representing different attitudes toward compensation, conjunction,
disjunction, and interaction among criteria. From this perspective, the
choice of aggregation operator is not merely a numerical convenience but
a modeling decision that determines how favorable and unfavorable
evidence may compensate for one another.

A substantial body of literature has explored these properties in
traditional decision-making settings. Reviews by
Mardani et al.~\citep{Mardani2018} and
Csisz\'ar~\citep{Csiszar2021} highlight the widespread use of OWA-based
aggregation and related fuzzy operators across diverse application
domains. More recent studies have continued to investigate how
aggregation semantics influence consensus formation, robustness,
criterion interaction, and ranking behavior. Examples include
consensus-oriented aggregation in group decision making
\citep{MoralEtAl2017,TrilloEtAl2023}, robust Choquet-based decision
models \citep{Yang2022Tolerance,Liao2023Contextual}, and applications
employing geometric and OWA-type operators to capture alternative
decision attitudes \citep{Stroppiana2022Explainable,Alcantud2023Multi}.

Despite this methodological richness, most aggregation studies are
formulated within classical MCDM settings, where the aggregated inputs
represent multiple criteria, attributes, preferences, or expert
judgments associated with the same decision problem. The primary
objective is typically to identify the most appropriate operator for
combining these criteria or to develop new aggregation mechanisms with
desirable mathematical properties. Comparatively less attention has been
paid to architectures in which the aggregated inputs originate from
fundamentally different decision layers.

This distinction is particularly important in cooperative automated
decision-making systems. In architectures such as CADEMAS-ML, the
aggregation process does not combine multiple criteria describing an
alternative. Instead, it reconciles two conceptually distinct sources of
evidence: a predictive assessment generated by machine-learning models
and a contextual assessment encoding organizational priorities,
regulatory requirements, expert knowledge, or policy constraints. Under
this interpretation, the aggregation operator effectively determines the
extent to which predictive evidence may compensate for weak contextual
support, or conversely, whether contextual constraints can restrict
otherwise favorable predictions. The operator therefore acts as an
implicit governance mechanism rather than merely a criterion-combination
tool.

For this reason, the present work focuses on three representative
aggregation regimes exhibiting fundamentally different compensation
semantics. The weighted arithmetic mean represents a fully compensatory
integration strategy in which predictive and contextual information may
offset one another. The weighted geometric mean represents a partially
compensatory strategy that penalizes strong imbalances between both
sources of evidence. Finally, the minimum operator represents a
non-compensatory strategy in which the weakest component determines the
final assessment. Together, these operators span a spectrum ranging from
highly permissive to highly restrictive integration behavior while
remaining interpretable and widely understood within the aggregation
literature.

From the perspective of this study, the objective is therefore not to
propose a new aggregation operator nor to extend aggregation theory.
Instead, aggregation operators are treated as alternative decision
semantics within a cooperative automated decision-making architecture.
The key question is how different compensation regimes influence policy
compliance, prioritization behavior, robustness to predictive
perturbations, and ranking stability when predictive and contextual
information must jointly determine the final recommendation.

\subsection{Research Gap and Positioning}
\label{subsec:research-gap}
% R1-1, R2-2, R2-3
% Diferenciar explícitamente la contribución respecto a:
% (i) agregación clásica en MCDM,
% (ii) lógica difusa,
% (iii) DSS sensibles al contexto.
% Introducir las preguntas de investigación.

The reviewed literature reveals two largely independent research
streams. On the one hand, context-aware and cooperative decision-support
systems have developed increasingly sophisticated mechanisms for
incorporating contextual knowledge into predictive processes through
ontologies, symbolic rules, knowledge graphs, human oversight, and
governance-oriented architectures
\citep{TengTan2008,YaoKumar2012,CronholmGoebel2022,Fanfani2025Human}.
These studies demonstrate the importance of contextual information but
typically focus on knowledge representation, rule selection, system
architecture, or human--AI interaction. Once predictive and contextual
assessments have been generated, the mechanism responsible for their
numerical integration is often treated as an implementation detail.

On the other hand, aggregation theory and fuzzy MCDM have developed a
rich collection of operators with diverse compensation semantics,
including OWA operators, geometric means, t-norms, Choquet integrals,
and Sugeno integrals
\citep{Yager1988,Yager2005,Mardani2018,Csiszar2021}. This literature
provides extensive theoretical knowledge regarding compensation,
criterion interaction, robustness, consensus formation, and preference
aggregation. However, aggregation operators are generally studied as
mathematical constructs or as components of multicriteria decision
methods, where the aggregated inputs represent criteria, preferences, or
expert judgments associated with the same decision problem.

Consequently, a gap exists between context-aware decision-support
frameworks and aggregation theory. The former provides mechanisms for
representing contextual knowledge but rarely investigates how different
aggregation semantics affect the preservation of contextual constraints
after integration. The latter offers powerful aggregation mechanisms but
seldom studies their implications when heterogeneous decision signals,
such as predictive machine-learning outputs and contextual assessments,
must jointly determine operational recommendations.

From this perspective, the gap addressed in this work is not the absence
of aggregation operators, nor the lack of mathematical knowledge about
their properties. Rather, it is the lack of a systematic understanding
of aggregation as a governance mechanism within cooperative automated
decision-making architectures. In particular, relatively little
attention has been devoted to questions such as whether contextual
assessments can effectively act as veto mechanisms, how different
compensation regimes influence policy compliance, how sensitive
prioritization outcomes are to predictive overconfidence, or how ranking
stability is affected by uncertainty in contextual information.

These observations motivate the following research questions:

\begin{itemize}
\item \textit{RQ1. How do different aggregation semantics affect the interaction between predictive evidence and contextual assessments in a cooperative automated decision-making architecture?}
\item \textit{RQ2. Under what conditions can aggregation operators preserve or violate contextual constraints and veto-like policies?}
\item \textit{RQ3. How do alternative aggregation regimes influence prioritization robustness under predictive overconfidence and contextual uncertainty?}
\item \textit{RQ4. To what extent do different aggregation semantics alter the stability of the resulting rankings and intervention priorities?}
\end{itemize}

Unlike traditional fuzzy MCDM studies, the objective of the present work
is not to propose a new aggregation operator, derive new algebraic
results, or solve a classical multicriteria decision problem.
Likewise, unlike many context-aware DSS frameworks, the focus is not on
developing new mechanisms for representing contextual knowledge.
Instead, aggregation operators are interpreted as alternative policy
semantics within the CADEMAS architecture. By keeping both the
predictive subsystem and the contextual subsystem fixed, the study
isolates the aggregation mechanism as an independent design variable and
examines its implications for policy compliance, robustness, and
prioritization behavior. This positioning distinguishes the contribution
from both aggregation theory and context-aware decision-support
research.

\subsection{Chronological Comparison of Related Work}
\label{subsec:chronological-comparison}
% R2-4
% Tabla comparativa.
% Destacar:
% - veto preservation,
% - robustness analysis,
% - policy-compliant integration,
% - cooperative ADM.

Table~\ref{tab:related-works-comparison} summarizes representative
contributions discussed in the previous subsections and positions the
present work with respect to the existing literature. The comparison
focuses on four dimensions that emerge from the identified research
gap: (i) the explicit integration of predictive and contextual
information, (ii) the treatment of contextual constraints or veto-like
policies, (iii) the analysis of aggregation robustness, and (iv) the
study of aggregation within cooperative automated decision-making
architectures.
 
\begin{table}[t]
\caption{Chronological comparison of representative related work. ``Partial'' indicates that a contribution addresses the dimension only indirectly (e.g., through rules, consensus thresholds, or architectural design) without a dedicated analysis of predictive--contextual score fusion.}
\label{tab:related-works-comparison}
\renewcommand{\arraystretch}{1.35}\footnotesize
\begin{adjustwidth}{-\extralength}{0cm}
\begin{tabularx}{\fulllength}{>{\raggedright\arraybackslash}l|>{\centering\arraybackslash}p{1.1cm}|>{\raggedright\arraybackslash}X|C|C|C|C}
\toprule
\textbf{Reference} &
\textbf{Year} & 
\textbf{Main focus} &
\textbf{Predictive + contextual integration} &
\textbf{Policy / veto analysis} &
\textbf{Aggregation robustness} &
\textbf{Cooperative ADM} \\
\midrule
Teng and Tan~\cite{TengTan2008} &
2008 &
Context-aware cognitive agents &
\checkmark & Partial & No & Partial \\
\midrule
Yao and Kumar~\cite{YaoKumar2012} &
2012 &
Ontology-based clinical DSS &
\checkmark & Partial & No & No \\
\midrule
Merkert et~al.~\cite{MerkertEtAl2015} &
2015 &
ML in DSS survey &
Partial & No & No & No \\
\midrule
Moral et~al.~\cite{MoralEtAl2017} &
2017 &
Consensus aggregation &
No & No & Partial & No \\
\midrule
Mardani et~al.~\cite{Mardani2018} &
2018 &
Fuzzy MCDM review &
No & No & No & No \\
\midrule
Csisz\'ar~\cite{Csiszar2021} &
2021 &
OWA theory &
No & Partial & Partial & No \\
\midrule
Cronholm and G\"obel~\cite{CronholmGoebel2022} &
2022 &
Human-centred AI DSS &
\checkmark & Partial & No & Partial \\
\midrule
Yang et~al.~\cite{Yang2022Tolerance} &
2022 &
Robust Choquet-based MCDM &
No & Partial & \checkmark & No \\
\midrule
Liao et~al.~\cite{Liao2023Contextual} &
2023 &
Contextual Choquet models &
Partial & No & \checkmark & No \\
\midrule
Trillo et~al.~\cite{TrilloEtAl2023} &
2023 &
Consensus-based GDM &
No & Partial & No & No \\
\midrule
Fanfani et~al.~\cite{Fanfani2025Human} &
2025 &
Human-centred AI survey &
\checkmark & Partial & No & Partial \\
\midrule
\textbf{Present work} &
2026 &
\textbf{CADEMAS aggregation semantics} &
\textbf{\checkmark} & \textbf{\checkmark} & \textbf{\checkmark} & \textbf{\checkmark} \\
\bottomrule
\end{tabularx}
\end{adjustwidth}
\end{table}

The table illustrates that prior research has typically focused either
on context-aware decision-support architectures or on aggregation
methods in MCDM and fuzzy decision making. Comparatively little
attention has been devoted to aggregation as a governance mechanism that
simultaneously integrates predictive evidence, contextual knowledge,
policy constraints, and robustness considerations within a cooperative
ADM framework.

%%%%%%%%%%%%%%%%%%%%%%%%%%%%%%%%%%%%%%%%%%%%%%%%%%%%%%%%%%%%



% R1-1, R2-2
% Mostrar por qué CADEMAS es el entorno adecuado para estudiar
% semánticas de agregación.

% R1-1:
% La novedad no está en los operadores sino en estudiar la agregación
% como componente arquitectónico independiente dentro de CADEMAS.

% R2-2:
% Diferencia explícita respecto a MCDM clásico:
% aquí no se agregan criterios homogéneos sino evidencia predictiva y
% contexto.

% R2-3:
% Conecta directamente el research gap con el modelo formal.

% R2-6:
% Justifica por qué estudiar operadores alternativos es posible y
% relevante dentro de CADEMAS.

% R1-6:
% Prepara la interpretación de los experimentos como análisis sistémico
% de una arquitectura y no como mera verificación de propiedades
% algebraicas.
\section{CADEMAS and CADEMAS-ML}
\label{sect:cademas}

The Cooperative Automated Decision-Making System (CADEMAS) framework
was introduced by \citet{NovoaHernandez2026} as a general architecture
for context-aware decision support in which multiple heterogeneous
Automated Decision-Making (ADM) units cooperate to produce a unified
recommendation, score, or decision. Rather than prescribing a specific
decision model, CADEMAS provides a modular structure that separates the
generation of decision evidence from its contextual evaluation and final
integration.

This separation is particularly relevant for the problem studied in the
present work. Many decision-support systems combine predictive evidence
with contextual information, but these elements are often intertwined
within the same model or decision process. In contrast, CADEMAS
explicitly distinguishes three conceptually different functions:
(i) evidence generation, (ii) contextual assessment, and
(iii) decision integration. Consequently, the aggregation mechanism can
be analysed independently of both the predictive subsystem and the
contextual knowledge representation.

Formally, a CADEMAS instance is represented by the tuple

\begin{equation}
\mathcal{S}=
\langle I,U,O,K,T\rangle,
\end{equation}
where $I$ denotes the Input Manager, $U$ the set of ADM units,
$O$ the Decision Integration Layer, $K$ the Contextual Modulator, and
$T$ the universe of admissible data representations.

Given an input case, the ADM units independently generate decision
evidence from the available information. In parallel, the Contextual
Modulator evaluates the extent to which the same case satisfies
organizational priorities, expert knowledge, regulatory requirements,
or other contextual constraints. The outputs of both subsystems are then
combined by the Decision Integration Layer to produce the final
decision score.

% TODO: The resulting architecture is illustrated in Figure~\ref{fig:cademas_architecture}. A defining characteristic of CADEMAS is that predictive evidence and contextual assessments remain separate until the final integration stage. This design makes the aggregation mechanism an explicit architectural component rather than an implicit consequence of the predictive model or the contextual representation.

An important specialization of this framework is CADEMAS-ML
\citep{Godz2026}, in which the ADM units correspond to machine-learning
models. In this setting, multiple predictive models cooperate to produce
a consolidated predictive assessment, while contextual information is
represented through a fuzzy inference system inspired by
\citet{PerezCanedo2022} and \citet{PerezCanedo2024}. The contextual
subsystem translates organizational priorities and institutional
requirements into a normalized degree of contextual satisfaction.

To focus on the role of aggregation, the internal implementation of both
subsystems is abstracted in terms of two normalized scores. For each
decision case $i$, the predictive subsystem produces a predictive score

\begin{equation}
R_i \in [0,1],
\end{equation}
which summarizes the cooperative evidence generated by the predictive
ADM units. Depending on the application domain, $R_i$ may represent a
risk estimate, a confidence score, a probability of occurrence, or any
other normalized predictive assessment.

Independently, the contextual subsystem produces a context score

\begin{equation}
Q_i \in [0,1],
\end{equation}
which quantifies the degree to which the case satisfies the contextual
requirements encoded by the organization, institution, or application
domain.

It is worth noting that the restriction of both predictive and
contextual scores to the unit interval does not limit the generality of
the proposed formulation. Any continuous score defined over an interval
$[a,b]$ can be transformed into the normalized domain $[0,1]$ through an
affine mapping. This transformation preserves ordering relations and
therefore maintains the ranking semantics required in prioritization
tasks. Consequently, $R_i$ and $Q_i$ should be interpreted as normalized
semantic scores rather than measurements expressed in a physical unit.
The use of the unit interval allows the analysis of aggregation
properties independently of the original scale of the predictive and
contextual information. More importantly, the lower boundary \(Q_i=0\) represents a complete \textit{contextual violation}.

The final prioritization score is then obtained through an aggregation
operator

\begin{equation}
P_i=A(R_i,Q_i),
\end{equation}
where $A$ defines the integration semantics between predictive evidence
and contextual assessment. Candidate cases are subsequently ranked
according to $P_i$, and the corresponding Top-$K$ subset is selected
for intervention or resource allocation.

It is important to emphasize that the aggregation problem considered
here differs from traditional fuzzy MCDM settings. In classical MCDM,
aggregation operators combine multiple criteria describing different
attributes of the same alternative, typically with the objective of
capturing trade-offs among decision criteria. In CADEMAS-ML, however,
the operator combines heterogeneous sources of decision evidence that
play fundamentally different roles: a predictive assessment generated by
cooperating machine-learning models and a contextual assessment derived
from organizational, regulatory, or domain-specific knowledge.

Consequently, the aggregation operator does not merely determine the
relative importance of criteria. Instead, it governs how predictive
evidence interacts with contextual requirements and therefore acts as an
explicit policy-design component within the decision architecture. The
central question is not which operator produces the most desirable
aggregate score in a multicriteria sense, but rather how alternative
aggregation semantics affect policy compliance, contextual veto
preservation, robustness to predictive overconfidence, and ranking
stability under uncertainty.

Existing CADEMAS-ML applications employ a weighted linear combination of
predictive and contextual scores \citep{Godz2026}. However, the
CADEMAS framework itself does not prescribe any particular aggregation
operator \citep{NovoaHernandez2026}. Consequently, different aggregation
semantics can be introduced without modifying either the predictive
subsystem or the contextual subsystem. This modularity provides an
ideal setting for studying aggregation as a policy-design mechanism and
for analysing how different compensation regimes affect prioritization,
policy compliance, robustness, and ranking stability. The following
section introduces the aggregation operators considered in this study
and develops the theoretical properties that motivate the subsequent
experimental analysis.

% R1-2, R1-3, R1-4, R1-5
% Reestructurar toda la teoría alrededor de propiedades.

%\subsection{Structural Properties of Aggregation Operators}

% R1-5, R2-5
% Definir:
% - boundedness
% - continuity
% - monotonicity
% - zero absorption
% - compensability
% - veto preservation
%
% Explicitar supuestos.

\section{Aggregation Semantics and Theoretical Analysis}
\label{sect:aggregation-theory}

The aggregation mechanism in CADEMAS-ML determines how predictive
evidence and contextual knowledge are integrated into a unified
prioritization decision. Therefore, the aggregation operator should not
be interpreted as a mere mathematical component, but as an explicit
design choice that defines the interaction between machine-learning
evidence and contextual requirements.

This section provides the theoretical foundations required to answer the
research questions introduced in Section~\ref{subsec:research-gap},
particularly those related to the role of aggregation semantics in
cooperative automated decision-making. Specifically, we analyze how
different aggregation regimes modify the interaction between predictive
and contextual information (RQ1), under which conditions contextual
constraints and veto-like policies are preserved or violated (RQ2), how
aggregation semantics influence robustness to predictive
overconfidence and contextual uncertainty (RQ3), and how they affect the
stability of the resulting prioritization rankings (RQ4).

The analysis does not aim to introduce a new aggregation operator or to
extend classical aggregation theory. Instead, it studies existing
aggregation paradigms as alternative policy semantics within the
CADEMAS-ML architecture. First, we formalize the structural properties
that characterize aggregation behavior, including boundedness,
monotonicity, continuity, zero absorption, contextual veto preservation,
and compensability. Then, we analyze three representative aggregation
regimes with different decision semantics: fully compensatory linear
aggregation, non-linear sub-compensatory aggregation through the weighted
geometric mean, and non-compensatory minimum aggregation. Finally, we
derive theoretical conditions describing contextual veto preservation,
rank reversal under compensation, and sensitivity to perturbations in
predictive and contextual assessments.

These results provide a formal basis for interpreting aggregation
operators as policy-design components in cooperative automated
decision-making systems and motivate the experimental evaluation of
their practical consequences in the following sections.


\subsection{Structural Properties of Aggregation Operators}
\label{subsec:structural-properties}

Without loss of generality, all scores are assumed to be normalized
within the unit interval. Any bounded continuous scoring scale
$[a,b]$ can be transformed into this representation through a monotonic
affine transformation, preserving the ordering relations and the
aggregation semantics considered in this work.

Let
\begin{equation}
A:[0,1]^2 \rightarrow [0,1]
\end{equation}
denote an aggregation operator that combines the predictive score
$R_i$ and the contextual score $Q_i$ into the final prioritization
score
\begin{equation}
P_i=A(R_i,Q_i).
\end{equation}

Throughout this work, aggregation operators are interpreted as
integration policies governing the interaction between predictive
evidence and contextual knowledge within the CADEMAS-ML architecture.
Rather than focusing on specific functional forms, we first introduce a
set of structural properties that characterize different aggregation
semantics.

\begin{Definition}[Boundedness]
An aggregation operator is bounded if
\begin{equation}
0 \leq A(r,q) \leq 1,
\qquad \forall (r,q)\in[0,1]^2.
\end{equation}
\end{Definition}
Boundedness guarantees that the resulting prioritization score remains
within the same normalized scale as the predictive and contextual
assessments.

\begin{Definition}[Monotonicity]
An aggregation operator is monotonic if

\begin{equation}
r_1 \leq r_2,\;
q_1 \leq q_2
\quad \Longrightarrow \quad
A(r_1,q_1)\leq A(r_2,q_2).
\end{equation}

\end{Definition}

Monotonicity ensures that improving either predictive evidence or
contextual satisfaction cannot decrease the final prioritization score.

\begin{Definition}[Continuity]
An aggregation operator is continuous if small perturbations in its
inputs produce arbitrarily small changes in the aggregated score.
\end{Definition}

Continuity is particularly relevant when predictive and contextual
scores are estimated from noisy data, since it prevents abrupt changes
in prioritization due to infinitesimal variations of the inputs.

\begin{Definition}[Zero Absorption]
An aggregation operator satisfies the zero-absorption property if

\begin{equation}
A(r,0)=0,
\qquad \forall r\in[0,1].
\end{equation}

\end{Definition}

Under zero absorption, a null contextual assessment completely
suppresses the final prioritization score regardless of the predictive
evidence available.

\begin{Definition}[Contextual Veto Preservation]
An aggregation operator preserves contextual vetoes if

\begin{equation}
Q_i=0
\quad \Longrightarrow \quad
P_i=0.
\end{equation}

\end{Definition}

Within the CADEMAS-ML setting, contextual veto preservation means that
contextual requirements cannot be overridden by predictive evidence.
Observe that zero absorption is a sufficient condition for contextual
veto preservation.

\begin{Definition}[Compensability]
An aggregation operator is compensatory if a reduction in one input can
be offset by a sufficiently large increase in the other input while
maintaining or improving the aggregated score.
\end{Definition}

Compensability determines the extent to which strong predictive evidence
may compensate for weak contextual support, or vice versa. At one
extreme, fully compensatory operators permit unrestricted trade-offs
between prediction and context. At the opposite extreme,
non-compensatory operators prevent such trade-offs and therefore behave
as hard contextual constraints.

These properties provide a theoretical basis for distinguishing
alternative integration semantics within CADEMAS-ML. In particular,
zero absorption and compensability play a central role in determining
whether contextual requirements behave as soft preferences, partial
constraints, or strict vetoes. The following subsection introduces the
specific aggregation operators considered in this work and analyzes how
they instantiate these structural properties.


\subsection{Aggregation Operators Considered}
\label{subsec:aggregation-operators}
% R1-4, R1-12, R2-6
% Linear.
% Weighted geometric.
% Minimum.
%
% Tratar cuidadosamente lambda=0 y lambda=1.

The CADEMAS framework does not prescribe a specific mechanism for
combining predictive and contextual assessments
\citep{NovoaHernandez2026}. Consequently, different aggregation
operators can be incorporated without modifying either the predictive
subsystem or the contextual subsystem.

To investigate how aggregation semantics affect prioritization
behavior, policy compliance, and robustness, this work considers three
representative operators that span different levels of compensability
between predictive evidence and contextual requirements.

\begin{enumerate}

\item \textit{Linear aggregation.}
The linear operator corresponds to the integration mechanism adopted in
existing CADEMAS-ML implementations \citep{Godz2026}:

\begin{equation}
A_L(R_i,Q_i)
=
\lambda R_i + (1-\lambda)Q_i,
\qquad \lambda\in[0,1].
\end{equation}

The parameter $\lambda$ controls the relative importance assigned to
predictive and contextual information. This operator is fully
compensatory because reductions in one component can always be offset by
sufficient increases in the other component.

\item \textit{Weighted geometric aggregation.}
As an intermediate alternative, we consider the weighted geometric
mean:

\begin{equation}
A_G(R_i,Q_i)
=
R_i^{\lambda}
Q_i^{1-\lambda},
\qquad \lambda\in(0,1).
\end{equation}

Unlike the linear operator, the weighted geometric mean penalizes
imbalanced evaluations and therefore reduces the degree of compensation
between predictive and contextual information. Nevertheless, partial
compensation remains possible whenever both inputs are strictly
positive.

\item \textit{Minimum aggregation.}
Finally, we consider the minimum operator

\begin{equation}
A_M(R_i,Q_i)
=
\min(R_i,Q_i).
\end{equation}

This operator represents a non-compensatory integration regime in which
the final prioritization score cannot exceed the weakest component.
Consequently, predictive evidence cannot compensate for insufficient
contextual support.

\end{enumerate}

Table~\ref{tab:operator-properties} summarizes the structural
properties of the three operators according to the definitions
introduced in the previous subsection.

\begin{table}[H]
\caption{Structural properties of the aggregation operators considered in this work.}
\label{tab:operator-properties}
\centering
\renewcommand{\arraystretch}{1.3}
\begin{tabular}{lccc}
\toprule
Property & Linear & Geometric & Minimum \\
\midrule
Boundedness & Yes & Yes & Yes \\
Monotonicity & Yes & Yes & Yes \\
Continuity & Yes & Yes & Yes \\
Zero absorption & No & Yes & Yes \\
Contextual veto preservation & No & Yes & Yes \\
Compensability & Full & Partial & No \\
\bottomrule
\end{tabular}
\end{table}

These operators were selected because they represent three distinct
integration semantics within CADEMAS-ML: unrestricted compensation,
partial compensation, and strict contextual constraint. The following
subsections analyze the theoretical consequences of these semantics and
their implications for prioritization behavior.

% R1-2, R1-3, R2-1
% Teorema general basado en zero absorption.
%
% Luego mostrar que minimum y geometric lo satisfacen.
\subsection{Conditions for Contextual Veto Preservation}
\label{subsec:veto-preservation}

Within CADEMAS-ML, contextual information may represent different
types of requirements, including organizational priorities, regulatory
constraints, safety conditions, or expert-defined restrictions. In some
applications, a contextual assessment should not behave as a soft
preference but rather as a hard constraint. This situation occurs when a
complete contextual violation must prevent an alternative from being
prioritized, regardless of its predictive evidence.

This behavior can be characterized through the zero-absorption property
introduced in Section~\ref{subsec:structural-properties}. An aggregation
operator provides contextual veto preservation if

\begin{equation}
A(r,0)=0,\qquad \forall r\in[0,1].
\end{equation}

Therefore, the question of whether CADEMAS-ML can enforce contextual
vetoes reduces to determining whether the selected aggregation operator
satisfies zero absorption.

\begin{Theorem}[Contextual veto preservation]
\label{thm:veto-preservation}
Let $A:[0,1]^2\rightarrow[0,1]$ be an aggregation operator used in
CADEMAS-ML. If $A$ satisfies zero absorption, then every decision case
$i$ with a null contextual assessment,
\begin{equation}
Q_i=0,
\end{equation}
receives a null prioritization score,
\begin{equation}
P_i=A(R_i,Q_i)=0,
\end{equation}
independently of its predictive score $R_i$.
\end{Theorem}

\begin{proof}
The result follows directly from the definition of zero absorption.
Since $Q_i=0$, we have $P_i=A(R_i,Q_i)=A(R_i,0)$.
Because $A$ satisfies zero absorption, $A(R_i,0)=0$, for every $R_i\in[0,1]$. Therefore, $P_i=0$.
Hence, a complete contextual violation completely suppresses the final
prioritization score regardless of the predictive evidence.
\end{proof}

\begin{Corollary}[Veto preservation of the considered operators]
    \label{cor:veto-operators}
    Among the aggregation operators considered in this work, the minimum
    and weighted geometric operators preserve contextual vetoes, whereas
    the linear operator does not.
\end{Corollary}
\begin{proof}
    For the minimum operator, $A_M(r,0)=\min(r,0)=0$.
    Similarly, for the weighted geometric operator, $A_G(r,0)=r^\lambda 0^{1-\lambda}=0$, for $\lambda\in(0,1)$. In contrast, the linear operator yields $A_L(r,0)=\lambda r$, which is greater than zero whenever $r>0$ and $\lambda>0$. Therefore, predictive evidence can override a complete contextual violation under linear aggregation.
\end{proof}

The previous result establishes that contextual veto preservation is
not a consequence of the predictive model or the contextual knowledge
representation, but rather an emergent property of the integration
semantics. Consequently, two CADEMAS-ML systems with identical
predictive and contextual components may exhibit fundamentally different
policy behavior depending only on the selected aggregation operator.

This observation motivates the subsequent analysis of rank reversal
conditions and compensation thresholds, which characterize how much
contextual improvement is required to offset predictive advantages under
different aggregation regimes.



\subsection{Conditions for Rank Reversal under Compensation}
\label{subsec:rank-reversal-compensation}
% R1-2, R2-1
% Generalizar más allá del operador lineal.
% Introducir compensatory aggregation.

The previous section established that contextual veto preservation depends on
the existence of a zero-absorption property. However, many aggregation
operators used in decision-support systems follow a different philosophy:
they allow partial compensation between the aggregated sources of evidence.
In such cases, a sufficiently high value in one dimension may offset a
lower value in another dimension.

This property is common in compensatory aggregation schemes, including
weighted arithmetic means and several non-linear aggregation operators.
Although compensation provides flexibility when predictive evidence and
contextual information represent partially substitutable criteria, it also
introduces the possibility that the final prioritization ranking differs
from the ranking induced by predictive evidence alone.

To formalize this phenomenon, let $A:[0,1]^2\rightarrow[0,1]$ be a
continuous aggregation operator that is strictly increasing in both
arguments. Given two alternatives $i$ and $j$, assume that alternative
$i$ has stronger predictive evidence, but alternative $j$ has better
contextual alignment:
$R_i > R_j$ and $Q_i < Q_j$.
The final prioritization scores are then
$P_i=A(R_i,Q_i)$ and $P_j=A(R_j,Q_j)$.
A rank reversal occurs when $P_j>P_i$,
meaning that the contextual advantage of alternative $j$ compensates for
its predictive disadvantage.

\begin{Theorem}[General rank reversal condition under compensation]
    \label{thm:general-rank-reversal}
    Let $A:[0,1]^2\rightarrow[0,1]$ be a continuous aggregation operator strictly increasing in both arguments. For two alternatives $i$ and $j$ satisfying predictive dominance $R_i>R_j$, assume the operator provides sufficient compensatory capacity such that there exists a contextual assessment $q_{\max} \le 1$ satisfying $A(R_j, q_{\max}) > A(R_i, Q_i)$.
    
    Under these conditions, a rank reversal in favour of alternative $j$ occurs if and only if
    $Q_j > Q_i+\Phi_A(R_i,R_j,Q_i)$,
    where $\Phi_A$ is the minimum contextual improvement required to compensate for the predictive difference, implicitly defined by $A(R_j,Q_i+\Phi_A)=A(R_i,Q_i)$.
\end{Theorem}
    
\begin{proof}
    Fix $R_i>R_j$ and $Q_i$. Define the difference function
    $g(q)=A(R_j,q)-A(R_i,Q_i)$.
    Because $A$ is continuous and strictly increasing in its second argument, $g(q)$ is also continuous and strictly increasing.
    At $q=Q_i$, $g(Q_i)=A(R_j,Q_i)-A(R_i,Q_i)<0$ because $A$ is strictly increasing in its first argument and $R_j<R_i$.
    By hypothesis, there exists $q_{\max}$ such that $g(q_{\max}) > 0$. Therefore, by the Intermediate Value Theorem, there exists a unique contextual value $q^\ast \in (Q_i, q_{\max})$ such that $g(q^\ast)=0$, which implies $A(R_j,q^\ast)=A(R_i,Q_i)$.
    Defining the required threshold as $\Phi_A(R_i,R_j,Q_i)=q^\ast-Q_i$, we obtain $P_j>P_i$ (or $g(Q_j)>0$) if and only if $Q_j>q^\ast$, which is equivalent to $Q_j>Q_i+\Phi_A(R_i,R_j,Q_i)$.
\end{proof}

The theorem shows that rank reversal is not an artefact of a particular
aggregation formula, but rather a consequence of allowing compensation
between predictive evidence and contextual information. The specific
threshold $\Phi_A$ depends on the semantics of each aggregation operator:
operators with stronger compensation require smaller contextual
improvements to reverse a predictive advantage, whereas operators with
reduced compensation require larger contextual differences.

This distinction is particularly relevant in CADEMAS-ML because the two
aggregated components do not represent equivalent evaluation criteria.
The predictive score $R_i$ summarizes evidence generated by machine
learning models, whereas the contextual score $Q_i$ represents compliance
with organizational, regulatory, or domain-specific requirements.
Therefore, the choice of aggregation operator implicitly defines the
degree to which contextual requirements can modify or override
data-driven recommendations.

\subsection{Sensitivity to Predictive and Contextual Perturbations}
\label{subsec:sensitivity-analysis}
% R1-9, R2-5
% Nuevo resultado teórico.
% Analizar: |A(R,Q+e)-A(R,Q)|
% Justificación formal de experimentos de robustez.

The previous sections analyzed the structural properties of aggregation
operators in terms of compensation and contextual veto preservation.
However, in practical decision-support systems, both predictive scores
and contextual assessments are subject to uncertainty. Predictive
models may produce unstable estimates due to limited data, model
uncertainty, or distribution changes, whereas contextual assessments may
vary because of incomplete knowledge, expert disagreement, or evolving
organizational requirements.

Therefore, an additional relevant property of an aggregation mechanism
is its sensitivity to perturbations in its input dimensions. In the
context of CADEMAS-ML, this property determines how strongly changes in
predictive evidence or contextual requirements affect the final
prioritization score.

Let $A:[0,1]^2\rightarrow[0,1]$ be an aggregation operator. Consider a
perturbation $\epsilon$ applied to the contextual assessment while the
predictive score remains fixed. The resulting variation in the
prioritization score is
$\Delta_Q(\epsilon)=\bigl|A(R,Q+\epsilon)-A(R,Q)\bigr|$.
Similarly, a perturbation in the predictive assessment produces
$\Delta_R(\epsilon)=\bigl|A(R+\epsilon,Q)-A(R,Q)\bigr|$.
These quantities measure the local influence of each information source
on the final decision score. Operators that strongly amplify small
perturbations may lead to unstable rankings, whereas operators with
bounded sensitivity provide greater robustness against uncertainty in
predictive or contextual information.

\begin{Theorem}[Lipschitz sensitivity of aggregation operators]
\label{thm:lipschitz-sensitivity}
Let $A:[0,1]^2\rightarrow[0,1]$ be an aggregation operator continuously
differentiable in the region containing
the line segment between $(R,Q)$ and $(R+\epsilon_R,Q+\epsilon_Q)$.
If there exist finite constants $L_R$ and $L_Q$ such that
\begin{equation*}
\bigl|\partial A/\partial R\bigr|\leq L_R \;\;\;\;\text{and}\;\;\;\;\;\; 
\bigl|\partial A/\partial Q\bigr|\leq L_Q,
\end{equation*}
then, for any perturbations $\epsilon_R$ and $\epsilon_Q$,
\begin{equation*}
\bigl|A(R+\epsilon_R,Q+\epsilon_Q)-A(R,Q)\bigr|
\leq L_R|\epsilon_R|+L_Q|\epsilon_Q|.
\end{equation*}
\end{Theorem}

\begin{proof}
By the Mean Value Theorem, there exists a point
$(\xi_R,\xi_Q)$ between $(R,Q)$ and
$(R+\epsilon_R,Q+\epsilon_Q)$ such that
\begin{equation*}
A(R+\epsilon_R,Q+\epsilon_Q)-A(R,Q)
=
(\partial A/\partial R)(\xi_R,\xi_Q)\epsilon_R
+
(\partial A/\partial Q)(\xi_R,\xi_Q)\epsilon_Q.
\end{equation*}

Taking absolute values and applying the derivative bounds yields
\begin{equation*}
\bigl|A(R+\epsilon_R,Q+\epsilon_Q)-A(R,Q)\bigr|
\leq L_R|\epsilon_R|+L_Q|\epsilon_Q|.
\end{equation*}
\end{proof}

The previous theorem establishes a general robustness condition for
aggregation mechanisms. The constants $L_R$ and $L_Q$ quantify the
maximum amplification of uncertainty from each decision component.
Consequently, aggregation operators with smaller sensitivity constants
produce more stable prioritization scores under uncertain inputs.

For the linear aggregation operator,
$A_L(R,Q)=\lambda R+(1-\lambda)Q$,
the sensitivity coefficients are constant:
$L_R=\lambda$ and $L_Q=1-\lambda$.
Therefore, the influence of predictive and contextual uncertainty is
uniform throughout the decision space. This property explains the fully
compensatory behaviour of the linear operator: a perturbation in one
dimension has exactly the same marginal effect regardless of the
current values of prediction and context.

For the weighted geometric operator,
$A_G(R,Q)=R^\lambda Q^{1-\lambda}$,
the local sensitivities are
$\partial A_G/\partial R=\lambda R^{\lambda-1}Q^{1-\lambda}$
and
$\partial A_G/\partial Q=(1-\lambda)R^\lambda Q^{-\lambda}$.
Unlike the linear operator, sensitivity is dependent on the current
state of the decision case. In particular, when $Q$ approaches zero,
contextual perturbations may have a stronger relative influence on the
final score. This behaviour reflects the non-linear penalization of
contextual incompatibility and provides a mathematical explanation for
the stronger robustness of non-compensatory and sub-compensatory aggregation mechanisms.

For the minimum operator,
$A_{\min}(R,Q)=\min(R,Q)$,
the final score is completely determined by the limiting component.
Consequently, the operator exhibits zero compensation: improving the
non-limiting input has no effect until it becomes the active constraint.
This property provides maximal preservation of contextual restrictions
but may reduce responsiveness when contextual assessments are uncertain. This property can be formilzed in the following corollary:

\begin{Corollary}[Minimum operator stability under inactive-component perturbations]
    \label{cor:min-noise}
    Consider a decision case evaluated under the minimum operator $A_M(R,Q) = \min(R,Q)$. If the predictive score acts as the strict limiting component ($R < Q$), the operator completely absorbs any contextual perturbation $\varepsilon$ as long as the ordinal relationship is maintained (i.e., $R \le Q + \varepsilon$). Under these conditions, the sensitivity is zero:
    $$ \bigl| A_M(R, Q+\varepsilon) - A_M(R, Q) \bigr| = |R - R| = 0. $$
\end{Corollary}
\begin{proof}
    Let the initial state of the decision case be defined by $R$ and $Q$. By hypothesis, the predictive score is the strict limiting component, meaning $R < Q$. By the definition of the minimum operator, the initial aggregated score is:
    $$ A_M(R, Q) = \min(R, Q) = R. $$
    
    Now, consider a perturbed contextual assessment $Q' = Q + \varepsilon$. We are given the condition that the perturbation $\varepsilon$ (whether positive or negative) is bounded such that the ordinal relationship is maintained, i.e., $R \le Q + \varepsilon$. Evaluating the minimum operator under this perturbed state yields:
    $$ A_M(R, Q+\varepsilon) = \min(R, Q+\varepsilon) = R. $$
    
    Calculating the absolute difference between the perturbed and initial states gives the sensitivity:
    $$ \bigl| A_M(R, Q+\varepsilon) - A_M(R, Q) \bigr| = | R - R | = 0. $$
    Thus, the perturbation is completely absorbed, and the final prioritization score remains unaffected.
\end{proof}

This result illustrates the asymmetric robustness of non-compensatory operators: they are insensitive to perturbations of the non-limiting component, but highly sensitive when perturbations modify the active constraint.

These analytical differences motivate the robustness experiments
conducted in the following sections. Rather than evaluating aggregation
operators only according to final ranking performance, we analyse how
their semantic properties influence stability under perturbations of
predictive evidence and contextual assessments.

%%%%%%%%%%%%%%%%%%%%%%%%%%%%%%%%%%%%%%%%%%%%%%%%%%%%%%%%%%%%

\section{Experimental Methodology}
\label{sect:experimental-methodology}
% R1-6, R1-7, R1-8, R2-7
% Separar claramente metodología de resultados.


The theoretical analysis established that aggregation operators induce
different decision semantics within CADEMAS-ML, affecting compensation,
contextual constraint preservation, and sensitivity to uncertainty.
However, the practical consequences of these properties emerge at the
system level, where multiple decision cases compete for prioritization
and intervention.

This section describes the experimental methodology designed to analyze
these effects. The experiments do not aim to verify the theoretical
results, but rather to investigate how different aggregation semantics
translate into observable decision-making behaviors under realistic
population-level scenarios. In particular, we evaluate their impact on
policy compliance, robustness against predictive overconfidence,
sensitivity to contextual uncertainty, and ranking stability.

The experimental design follows a Monte Carlo simulation framework in
which synthetic decision populations are generated under controlled
conditions. This approach allows the independent manipulation of
predictive evidence, contextual assessments, uncertainty levels, and
aggregation regimes while keeping the underlying decision architecture
unchanged. Consequently, differences observed across scenarios can be
attributed to the integration semantics rather than to variations in the
predictive or contextual subsystems.

\subsection{Objectives and Experimental Logic}
\label{subsec:experimental-objectives}
% R1-6
% Explicar que Monte Carlo NO verifica la teoría.
%
% Evalúa efectos poblacionales emergentes:
% - Top-K,
% - robustness,
% - policy compliance.

The experimental evaluation is designed to answer how alternative
aggregation semantics affect the operational behavior of CADEMAS-ML when
applied to collections of competing decision cases. While the theoretical
analysis characterizes individual properties of aggregation operators,
the experimental analysis examines their emergent consequences when
scores are used to prioritize a population of alternatives.

Specifically, the experiments focus on four aspects. First, we analyze
policy compliance by measuring whether alternatives violating contextual
requirements can still reach the highest-priority positions under
different aggregation regimes. Second, we study the effects of
predictive overconfidence, where machine-learning evidence is assumed to
be systematically optimistic despite poor contextual alignment. Third,
we evaluate robustness to contextual uncertainty by introducing
controlled perturbations in contextual assessments. Finally, we analyze
the stability of the resulting prioritization decisions through ranking
and Top-$K$ consistency measures.

For each simulated scenario, a population of decision cases is generated,
where each case $i$ is characterized by a predictive score
$R_i\in[0,1]$ and a contextual score $Q_i\in[0,1]$. A given aggregation
operator $A$ produces the final prioritization score
\[
P_i=A(R_i,Q_i),
\]
which determines the ranking of alternatives and the selection of the
Top-$K$ candidates. The same population and perturbation conditions are
evaluated under all aggregation regimes, allowing a controlled
comparison of their decision semantics. The aggregation parameter $\lambda$ is varied according to the specific operator definition, allowing the influence of predictive and contextual information to be systematically analyzed while preserving the normalized decision scale.


% R1-7
% Distribuciones.
% Parámetros.
% Seeds.
\subsection{Synthetic Population Generation}
\label{subsec:synthetic-population}

The objective of the synthetic experiments is not to reproduce a
particular application domain, but to create controlled decision
environments in which the effects of different aggregation semantics can
be isolated. Therefore, synthetic populations are generated by directly
sampling the two components that define each CADEMAS-ML decision case:
the predictive assessment and the contextual assessment.

Each simulated population consists of $N$ independent decision cases.
For each case $i$, a predictive score $R_i\in[0,1]$ and a contextual
assessment $Q_i\in[0,1]$ are generated. These values represent the outputs
of the predictive and contextual subsystems, respectively, while the
internal mechanisms used to obtain them are abstracted away. This design
allows the aggregation stage to be evaluated independently from any
specific machine-learning model or knowledge representation approach.

To generate realistic score distributions, predictive and contextual
values are sampled from Beta distributions:
\begin{equation}
R_i\sim\mathrm{Beta}(\alpha_R,\beta_R),
\qquad
Q_i\sim\mathrm{Beta}(\alpha_Q,\beta_Q),
\end{equation}
which provide flexible control over concentration, dispersion, and
asymmetry within the unit interval. The Beta distribution is particularly
suitable for normalized confidence, risk, and satisfaction scores because
it naturally represents bounded continuous assessments.

Different experimental scenarios are obtained by modifying the parameters
of these distributions and by introducing controlled perturbations or
dependencies between predictive and contextual assessments when required.
In particular, three main scenarios are considered:

\begin{enumerate}

\item \textit{Baseline decision environment.}
Predictive and contextual assessments are generated independently,
representing situations where machine-learning evidence and contextual
knowledge provide complementary but unrelated information.

\item \textit{Predictive overconfidence scenario.}
Predictive overconfidence is modeled as a systematic upward bias in the
predictive assessment. Specifically, the original predictive score is
modified as

\begin{equation}
R_i'=\mathrm{clip}(R_i+\delta,0,1),
\end{equation}

where $\delta\geq0$ represents the magnitude of the artificial
overconfidence introduced into the predictive subsystem and
$\mathrm{clip}(\cdot)$ ensures that the perturbed values remain within the
normalized domain. This scenario represents situations in which machine
learning models assign excessively high predictive scores despite limited
contextual support.

\item \textit{Contextual uncertainty scenario.}
Contextual uncertainty is simulated by introducing stochastic
perturbations into the contextual assessments:

\begin{equation}
Q_i'=\mathrm{clip}(Q_i+\epsilon_i,0,1),
\end{equation}

where

\begin{equation}
\epsilon_i\sim\mathcal{N}(0,\sigma_Q^2)
\end{equation}

represents contextual uncertainty and $\sigma_Q$ controls its intensity.
The clipping operation guarantees that perturbed contextual assessments
remain within the valid normalized range. This formulation models
uncertainty arising from incomplete contextual knowledge, expert
disagreement, or variations in organizational requirements.

\end{enumerate}

For each scenario, identical synthetic populations are evaluated using
all aggregation operators. This paired experimental design ensures that
observed differences in prioritization, policy compliance, and ranking
stability are attributable to the aggregation semantics rather than to
variations in the generated decision cases.

All simulations are repeated using multiple random seeds to account for
sampling variability. The number of generated cases, distribution
parameters, perturbation levels, and random seeds are reported in the
following subsection to ensure experimental reproducibility.
\subsection{Monte Carlo Protocol}

% R1-7
% Aumentar replicaciones.
%
% Explicar CI.
%
% Reproducibilidad.

\subsection{Evaluation Metrics}

% R1-11
% Policy violation.
% Average predictive score.
% Manejar denominadores vacíos.

% R1-8
% Kendall tau.
% Jaccard Top-K.

% R1-4
% Tratamiento de empates.

%%%%%%%%%%%%%%%%%%%%%%%%%%%%%%%%%%%%%%%%%%%%%%%%%%%%%%%%%%%%

\section{Simulation Results}

% R1-6, R1-8, R2-7
% Aquí van únicamente resultados.

\subsection{Policy Compliance under Contextual Vetoes}

% R1-6
% Mostrar consecuencias sistémicas.

\subsection{Effects of Predictive Overconfidence}

% R1-6
% Definir claramente overconfidence artificial.

\subsection{Robustness to Contextual Uncertainty}

% R1-8, R1-9, R2-7
% Variar lambda.
% Variar ruido.
% Variar población.

\subsection{Sensitivity Analysis}

% R1-8, R2-7
% Consolidar todos los análisis de sensibilidad.

%%%%%%%%%%%%%%%%%%%%%%%%%%%%%%%%%%%%%%%%%%%%%%%%%%%%%%%%%%%%

\section{Employee Attrition Case Study}

% R1-10
% Presentarlo como ilustración.
% No como validación independiente.

\subsection{Original CADEMAS-ML Application}

% Explicar origen de scores.

\subsection{Experimental Setup}

% R1-10
% Top-K.
% Empates.
% Context scores.
% Procedimiento.

\subsection{Results and Interpretation}

% R1-10
% Mostrar estabilidad para distintos K.

%%%%%%%%%%%%%%%%%%%%%%%%%%%%%%%%%%%%%%%%%%%%%%%%%%%%%%%%%%%%

\section{Discussion}

% R1-1, R1-8, R2-2

\subsection{Implications for Cooperative Automated Decision Making}

% Discusión conceptual.

\subsection{Relation to Fuzzy MCDM and Aggregation Theory}

% R1-1, R2-2
% Subsection MUY IMPORTANTE.
%
% Explicar explícitamente:
% - qué NO aporta el paper;
% - qué sí aporta.
%
% Distinguir:
% aggregation theory,
% fuzzy MCDM,
% CADEMAS.

\subsection{Limitations}

% R1-8, R1-10
% Limitar alcance de conclusiones.

%%%%%%%%%%%%%%%%%%%%%%%%%%%%%%%%%%%%%%%%%%%%%%%%%%%%%%%%%%%%

\section{Conclusions and Future Work}

% R1-8
% Moderar afirmaciones.
%
% No decir:
% "nonlinear operators are superior".
%
% Decir:
% "under the considered conditions".

% Futuro:
% OWA,
% Choquet,
% learned aggregation,
% dynamic policies.

%%%%%%%%%%%%%%%%%%%%%%%%%%%%%%%%%%%%%%%%%%
\authorcontributions{To be completed.}

\funding{This research was supported by the project ``Study, Analysis and Evaluation of Cooperative Automated Decision-Making Systems (CADEMAS)'', PID2023-146575NB-I00, funded by MCIU/AEI/I0.13039/501100011033 and by FSE+.}

\institutionalreview{Not applicable.}

\informedconsent{Not applicable.}

\dataavailability{To be completed.}

\acknowledgments{Mariia Godz has a pre-doctoral research contract associated with the project ``Study, Analysis and Evaluation of Cooperative Automated Decision-Making Systems (CADEMAS)'', PID2023-146575NB-I00, funded by MCIU/AEI/I0.13039/501100011033 and by FSE+.}

\conflictsofinterest{The authors declare no conflicts of interest.}

%%%%%%%%%%%%%%%%%%%%%%%%%%%%%%%%%%%%%%%%%%
\begin{adjustwidth}{-\extralength}{0cm}
\reftitle{References}
%\bibliography{biblio}

\begin{thebibliography}{999}

\bibitem[Teng and Tan(2008)]{TengTan2008}
Teng, T.H.; Tan, A.H.
\newblock Cognitive Agents Integrating Rules and Reinforcement Learning for Context-Aware Decision Support.
\newblock In Proceedings of the Proceedings of the 2008 IEEE/WIC/ACM International Conference on Web Intelligence and Intelligent Agent Technology,  2008, pp. 318--321.
\newblock {\url{https://doi.org/10.1109/WIIAT.2008.163}}.

\bibitem[Yao and Kumar(2012)]{YaoKumar2012}
Yao, W.; Kumar, A.
\newblock Integrating Clinical Pathways into {CDSS} Using Context and Rules: A Case Study in Heart Disease.
\newblock In Proceedings of the Proceedings of the 2nd ACM SIGHIT International Health Informatics Symposium,  2012, pp. 267--276.
\newblock {\url{https://doi.org/10.1145/2110363.2110431}}.

\bibitem[Abu-Rasheed et~al.(2024)Abu-Rasheed, Abdulsalam, Weber, and Fathi]{AbuRasheed2024}
Abu-Rasheed, H.; Abdulsalam, M.; Weber, C.; Fathi, M.
\newblock Supporting Student Decisions on Learning Recommendations: An {LLM}-Based Chatbot with Knowledge Graph Contextualization for Conversational Explainability and Mentoring.
\newblock {\em Frontiers in Artificial Intelligence} {\bf 2024}.
\newblock {\url{https://doi.org/10.35542/osf.io/ervym}}.

\bibitem[Casado and Fernandez(2025)]{Casado2025Governance}
Casado, D.F.; Fernandez, R.P.
\newblock A Governance-Aware Methodological Framework for Cognitive {AI} Integration in Marine {CAD}: A Conceptual Pipe Routing Use Case.
\newblock {\em Journal of Marine Science and Engineering} {\bf 2025}, {\em 14},~1352.
\newblock {\url{https://doi.org/10.3390/jmse14011352}}.

\bibitem[Cronholm and G{\"o}bel(2022)]{CronholmGoebel2022}
Cronholm, S.; G{\"o}bel, H.
\newblock Design Principles for Human-Centred {AI}.
\newblock In Proceedings of the ECIS 2022 Research Papers,  2022.
\newblock Research paper 32.

\bibitem[Fanfani et~al.(2025)Fanfani, Palesi, and Nesi]{Fanfani2025Human}
Fanfani, M.; Palesi, L.A.I.; Nesi, P.
\newblock Human-Centered {AI} for Decision Support Systems: A Systematic Review of Application Domains, Architecture Designs, Current Trends and Future Directions.
\newblock {\em Big Data and Cognitive Computing} {\bf 2025}, {\em 10},~186.
\newblock {\url{https://doi.org/10.3390/bdcc10060186}}.

\bibitem[Smirnov et~al.(2021)Smirnov, Levashova, Ponomarev, and Shilov]{Smirnov2021Methodology}
Smirnov, A.V.; Levashova, T.; Ponomarev, A.; Shilov, N.
\newblock Methodology for Multi-Aspect Ontology Development: Ontology for Decision Support Based on Human-Machine Collective Intelligence.
\newblock {\em IEEE Access} {\bf 2021}, {\em 9},~130721--130739.
\newblock {\url{https://doi.org/10.1109/ACCESS.2021.3116870}}.

\bibitem[Marot et~al.(2022)Marot, Kelly, Naglic, Barbesant, and Cremer]{Marot2022Perspectives}
Marot, A.; Kelly, A.; Naglic, M.; Barbesant, V.; Cremer, J.L.
\newblock Perspectives on Future Power System Control Centers for Energy Transition.
\newblock {\em Journal of Modern Power Systems and Clean Energy} {\bf 2022}, {\em 10},~328--344.
\newblock {\url{https://doi.org/10.35833/mpce.2021.000673}}.

\bibitem[Novoa-Hernández et~al.(2026)Novoa-Hernández, Pelta, Godz, Verdegay, and Buendia-Carrillo]{NovoaHernandez2026}
Novoa-Hernández, P.; Pelta, D.; Godz, M.; Verdegay, J.; Buendia-Carrillo, D.
\newblock {CADEMAS} – A Framework for Cooperative Automated Decision-Making Systems.
\newblock In Proceedings of the Proceedings of the 2026 IEEE Conference on Artificial Intelligence (CAI),  2026, pp. 756--761.
\newblock {\url{https://doi.org/10.1109/CAI68641.2026.11536392}}.

\bibitem[Godz et~al.(2026)Godz, Novoa-Hernández, and Pelta]{Godz2026}
Godz, M.; Novoa-Hernández, P.; Pelta, D.
\newblock A Cooperative and Context-Aware Approach for Employee Attrition Prevention. In {\em Information Processing and Management of Uncertainty in Knowledge-Based Systems}; Springer Nature Switzerland,  2026; pp. 203--217.
\newblock {\url{https://doi.org/10.1007/978-3-032-29000-7_15}}.

\bibitem[Yager(1988)]{Yager1988}
Yager, R.
\newblock On Ordered Weighted Averaging Aggregation Operators in Multicriteria Decisionmaking.
\newblock {\em IEEE Transactions on Systems, Man, and Cybernetics} {\bf 1988}, {\em 18},~183--190.
\newblock {\url{https://doi.org/10.1109/21.87068}}.

\bibitem[Yager(2005)]{Yager2005}
Yager, R.R.
\newblock Extending Multicriteria Decision Making by Mixing t-Norms and {OWA} Operators.
\newblock {\em International Journal of Intelligent Systems} {\bf 2005}, {\em 20},~453--474.
\newblock {\url{https://doi.org/10.1002/int.20075}}.

\bibitem[Mardani et~al.(2018)Mardani, Nilashi, Zavadskas, Awang, Zare, and Jamal]{Mardani2018}
Mardani, A.; Nilashi, M.; Zavadskas, E.K.; Awang, S.; Zare, H.; Jamal, N.M.
\newblock Decision making methods based on fuzzy aggregation operators: Three decades review from 1986 to 2017.
\newblock {\em International Journal of Information Technology \& Decision Making} {\bf 2018}, {\em 17},~391--466.
\newblock {\url{https://doi.org/10.1142/S021962201830001X}}.

\bibitem[Csisz{\'a}r(2021)]{Csiszar2021}
Csisz{\'a}r, O.
\newblock Ordered Weighted Averaging Operators: A Short Review.
\newblock {\em IEEE Systems, Man, and Cybernetics Magazine} {\bf 2021}, {\em 7},~4--12.
\newblock {\url{https://doi.org/10.1109/MSMC.2020.3036378}}.

\bibitem[del Moral et~al.(2017)del Moral, Chiclana, Tapia~Garc{\'i}a, and Herrera-Viedma]{MoralEtAl2017}
del Moral, M.J.; Chiclana, F.; Tapia~Garc{\'i}a, J.M.; Herrera-Viedma, E.
\newblock An Analysis on Consensus Measures in Group Decision Making.
\newblock In Proceedings of the Proceedings of the 2017 4th International Conference on Control, Decision and Information Technologies (CoDIT),  2017, pp. 0262--0267.
\newblock {\url{https://doi.org/10.1109/CoDIT.2017.8102605}}.

\bibitem[Trillo et~al.(2023)Trillo, Cabrerizo, del Moral, Morente-Molinera, Tapia, and Herrera-Viedma]{TrilloEtAl2023}
Trillo, J.R.; Cabrerizo, F.J.; del Moral, M.J.; Morente-Molinera, J.A.; Tapia, J.M.; Herrera-Viedma, E.
\newblock A Consensus-Based Multi-Criteria Group Decision-Making Method Based on an Aggregated Operator Customised by Experts.
\newblock In Proceedings of the Proceedings of the 2023 9th International Conference on Control, Decision and Information Technologies (CoDIT),  2023.
\newblock {\url{https://doi.org/10.1109/CoDIT58514.2023.10284209}}.

\bibitem[Yang et~al.(2022)Yang, Lin, Fu, and Luo]{Yang2022Tolerance}
Yang, Y.W.; Lin, J.K.; Fu, Y.; Luo, Y.
\newblock Tolerance Framework for Robust Group Multiple Criteria Decision Making.
\newblock {\em Expert Systems with Applications} {\bf 2022}, {\em 204},~117430.
\newblock {\url{https://doi.org/10.1016/j.eswa.2022.118208}}.

\bibitem[Liao et~al.(2023)Liao, Liao, and Zhang]{Liao2023Contextual}
Liao, Z.; Liao, H.; Zhang, X.
\newblock A Contextual {Choquet} Integral-Based Preference Learning Model Considering Both Criteria Interactions and the Compromise Effects of Decision-Makers.
\newblock {\em Expert Systems with Applications} {\bf 2023}, {\em 213},~118977.
\newblock {\url{https://doi.org/10.1016/j.eswa.2022.118977}}.

\bibitem[Stroppiana et~al.(2022)Stroppiana, Boschetti, Brivio, Manfron, and Antoninetti]{Stroppiana2022Explainable}
Stroppiana, D.; Boschetti, M.; Brivio, P.A.; Manfron, G.; Antoninetti, M.
\newblock Explainable Multi-Criteria Data-Driven Environmental Status Assessment from Remote Sensing.
\newblock In Proceedings of the Proceedings of the IEEE MELECON 2022,  2022.
\newblock {\url{https://doi.org/10.1109/MELECON53508.2022.9842919}}.

\bibitem[Alcantud(2023)]{Alcantud2023Multi}
Alcantud, J.C.R.
\newblock Multi-Attribute Group Decision-Making Based on Intuitionistic Fuzzy Aggregation Operators Defined by Weighted Geometric Means.
\newblock {\em Granular Computing} {\bf 2023}, {\em 8},~1857--1866.
\newblock {\url{https://doi.org/10.1007/s41066-023-00406-w}}.

\bibitem[Merkert et~al.(2015)Merkert, Mueller, and Hubl]{MerkertEtAl2015}
Merkert, J.; Mueller, M.; Hubl, M.
\newblock A Survey of the Application of Machine Learning in Decision Support Systems.
\newblock In Proceedings of the Proceedings of the 23rd European Conference on Information Systems (ECIS),  2015.
\newblock {\url{https://doi.org/10.18151/7217429}}.

\bibitem[Pérez-Cañedo et~al.(2023)Pérez-Cañedo, Porras, Pelta, and Verdegay]{PerezCanedo2022}
Pérez-Cañedo, B.; Porras, C.; Pelta, D.; Verdegay, J.
\newblock Modeling Contexts as Fuzzy Propositions in Optimization Problems.
\newblock {\em {IEEE} Transactions on Fuzzy Systems} {\bf 2023}, {\em 31},~1474--1483.
\newblock {\url{https://doi.org/10.1109/tfuzz.2022.3203786}}.

\bibitem[P{\'e}rez-Ca{\~n}edo et~al.(2024)P{\'e}rez-Ca{\~n}edo, Novoa-Hern{\'a}ndez, Porras, Pelta, and Verdegay]{PerezCanedo2024}
P{\'e}rez-Ca{\~n}edo, B.; Novoa-Hern{\'a}ndez, P.; Porras, C.; Pelta, D.; Verdegay, J.
\newblock Contextual analysis of solutions in a tourist trip design problem: A fuzzy logic-based approach.
\newblock {\em Applied Soft Computing} {\bf 2024}, {\em 154},~111351.
\newblock {\url{https://doi.org/10.1016/j.asoc.2024.111351}}.

\end{thebibliography}
\end{adjustwidth}
\PublishersNote{}

\end{document}
